% Self-contained technical note for internal project documentation.
% Rationale for the sequential-approximation strategy used in the
% eigenfrequency topology-optimization implementation, framed as a
% reinterpretation of the role of the local optimization subproblem in the
% nested inner/outer procedure of Du and Olhoff (2007).
%
% Notation follows Du & Olhoff (2007), Struct Multidisc Optim 34:91-110,
% and its erratum (s00158-007-0167-6).
\documentclass[11pt]{article}
\usepackage{amsmath,amssymb}
\usepackage{bm}
\usepackage[margin=1in]{geometry}

\newcommand{\rhov}{\bm{\rho}}
\newcommand{\Krho}{\mathbf{K}}
\newcommand{\Mrho}{\mathbf{M}}
\newcommand{\phiv}{\bm{\phi}}
\newcommand{\fv}{\mathbf{f}}
\newcommand{\NE}{N_{\mathrm{E}}}
\newcommand{\T}{^{\mathsf{T}}}

\title{Sequential Approximation Strategy for Eigenfrequency
Maximization:\\ Rationale for the Present Implementation}
\author{}
\date{}

\begin{document}
\maketitle

\section{The reference algorithm of Du and Olhoff}

We consider the maximization of the $n$-th natural eigenfrequency
$\omega_n$ (or of a frequency gap) of a freely vibrating continuum
discretized into $\NE$ finite elements with density design variables
$\rho_e$, $e=1,\dots,\NE$, subject to a volume constraint. With
$\lambda_j=\omega_j^2$, the state is governed by the generalized
eigenvalue problem
\begin{equation}
\Krho(\rhov)\,\phiv_j = \lambda_j\,\Mrho(\rhov)\,\phiv_j,
\qquad \phiv_j\T\Mrho\,\phiv_k=\delta_{jk}.
\label{eq:gevp}
\end{equation}
Following Du and Olhoff (2007), the problem is cast in \emph{bound}
(max--min) form: a scalar $\beta$ is maximized subject to
$\beta\le\lambda_j$ for the relevant modes, so that raising $\beta$
raises the lowest active eigenvalue. This formulation is essential
because the eigenvalue of interest may be \emph{multiple}: when
$\lambda_n=\dots=\lambda_{n+N-1}$ with multiplicity $N>1$, the
eigenvalues are only directionally differentiable, and the bound
variable absorbs the resulting non-smoothness.

The computational procedure (their Fig.~1) is a \emph{nested} scheme
with a main (outer) loop and an inner loop. Each outer iteration, at the
current design $\rhov$, performs three analysis steps: (i) solve
\eqref{eq:gevp} by finite element analysis and detect the multiplicity
$N$ of $\omega_n$ within a small relative tolerance; (ii) assemble the
\emph{generalized gradient vectors} $\fv_{sk}$, $s,k=n,\dots,n+N-1$,
whose components are
\begin{equation}
\big(\fv_{sk}\big)_e=
\phiv_s\T\!\left(\Krho'_{\rho_e}-\lambda_n\,\Mrho'_{\rho_e}\right)\!\phiv_k,
\label{eq:fsk}
\end{equation}
reducing for a simple eigenvalue ($N=1$) to the ordinary sensitivity
$\fv_{nn}=\nabla\lambda_n$; and (iii) freeze
$\Krho,\Mrho,\lambda_j,\phiv_j,N,\fv_{sk}$ and solve, by MMA, a
sub-optimization problem for the design \emph{increments} $\Delta\rho_e$.
In compact form this frozen subproblem reads
\begin{subequations}
\label{eq:sub}
\begin{align}
\max_{\beta,\,\Delta\rhov}\ & \beta \\
\text{s.t.}\quad
 \beta &\le \lambda_j+\fv_{jj}\T\Delta\rhov, & j&=n+N,
 \label{eq:sub-guard}\\
 \beta &\le \lambda_j+\Delta\lambda_j, & j&=n,\dots,n+N-1,
 \label{eq:sub-mult}\\
 0 &=\det\!\big[\fv_{sk}\T\Delta\rhov-\delta_{sk}\,\Delta\lambda\big],
 & s,k&=n,\dots,n+N-1, \label{eq:sub-det}\\
 &\textstyle\sum_{e}(\rho_e+\Delta\rho_e)V_e \le V^*, \\
 &\ \rho_{\min}\le \rho_e+\Delta\rho_e\le 1, & e&=1,\dots,\NE.
 \label{eq:sub-box}
\end{align}
\end{subequations}
Here the multiple-eigenvalue increments $\Delta\lambda_j$ are the roots of
the algebraic sub-eigenvalue problem \eqref{eq:sub-det} (with the erratum
correction placing $\Delta$ on the eigenvalue $\Delta\lambda$), and
\eqref{eq:sub-guard} is a linearized guard on the simple mode just above
the cluster. The inner loop iterates \eqref{eq:sub} to convergence in
$\Delta\rhov$; the design is then updated as
$\rho_e\!:=\!\rho_e+\Delta\rho_e$ and the outer loop terminates once
$\|\Delta\rhov\|<\varepsilon$. The inner loop exists precisely for the
multiple case: through \eqref{eq:sub-det} the $\Delta\lambda_j$ are
\emph{nonlinear} functions of $\Delta\rhov$, so the frozen subproblem is a
genuinely nonlinear program that MMA must resolve iteratively. For $N=1$,
\eqref{eq:sub-det} collapses to the scalar relation
$\Delta\lambda_n=\fv_{nn}\T\Delta\rhov$ and \eqref{eq:sub} becomes linear
in $\Delta\rhov$.

\section{The present implementation}

Our implementation retains the bound formulation and MMA but replaces the
nested procedure by a single-loop sequential scheme. After each solution
of \eqref{eq:gevp} we assemble a convex separable approximation of the
bound problem and take \emph{one} MMA step on the absolute design
variables $\rho_e$; the eigenproblem is then recomputed immediately at the
new design. We do not detect multiplicity and do not form the generalized
gradients \eqref{eq:fsk}: the bound variable is carried as a persistent
MMA design variable, and the several lowest modes $\omega_1,\dots,\omega_J$
are constrained simultaneously, $\beta\le\lambda_j$, $j=1,\dots,J$. The
step is regularized by a move limit, by density filtering, by a Heaviside
projection under a continuation schedule, and by an acceptance test that
rejects any step which reduces $\omega_1$ beyond a small tolerance.

The governing idea is a change of interpretation rather than of
objective. We treat each finite element eigensolution as furnishing a
local model of the eigenvalue landscape that is trustworthy only in a
neighbourhood of the current design, and we advance by a single
conservative step within that neighbourhood before rebuilding the model.
What this reflects about the role of the frozen subproblem
\eqref{eq:sub}---whether it is an optimization problem to be solved, or a
predictor of a search direction---is developed in the next section, and it
is there, rather than in the reduction to a single step, that the essential
difference from the reference algorithm resides.

\section{Mathematical interpretation}

\subsection{Local optimization versus local prediction}

The frozen subproblem \eqref{eq:sub} admits two mathematically legitimate
readings, and the difference between them---rather than any single
implementation detail---is what distinguishes the present scheme from the
reference algorithm.

Under the first reading, \eqref{eq:sub} is itself an optimization problem:
given the current physical model, one seeks the increment $\Delta\rhov$
that best improves the bound and solves this subproblem close to
convergence before the finite element model is refreshed. This is the
reading embodied by the nested inner loop of Du and Olhoff, and it follows
naturally from the mathematical concern that shaped their formulation. Their
central difficulty was the treatment of \emph{multiple} eigenvalues: at a
multiplicity the ordinary eigenvalue sensitivities cease to exist, and the
generalized gradients \eqref{eq:fsk} together with the coupled determinant
\eqref{eq:sub-det} were introduced precisely to furnish a well-defined local
model of the degenerate cluster. Once a local model of that richness
had been constructed, solving it accurately before rebuilding is a natural
design decision, since the subproblem then encodes genuine directional
information about the multiple eigenvalue that is worth extracting in full.
In this reading the governing question is \emph{how accurately the local
approximation should be solved}.

Under the second reading, adopted here, \eqref{eq:sub} is not an end in
itself but a \emph{local predictor} of a useful search direction. It
converts the frozen sensitivities into a descent direction; once a single
conservative step has been taken, the eigendata on which that direction was
built no longer describe the current design, and the proper response is to
rebuild the model rather than to refine the increment further. The governing
question here is therefore not how accurately the subproblem is solved but
\emph{how long the local approximation remains trustworthy}. The limitation
to a single MMA step reflects no deficiency of MMA as an optimizer: after one
conservative update the optimization problem has itself changed, since the
stiffness and mass matrices $\Krho,\Mrho$, the eigenvectors $\phiv_j$, and
the eigenvalue sensitivities now all correspond to a different design. The
decision to rebuild is thus prompted by the change in the underlying
physical model rather than by any shortcoming of the solver; and because a
fresh model is inexpensive to form here, keeping the approximation current
is preferred to solving an increasingly outdated one to high accuracy.

The two readings differ in where convergence is expected to arise. The
first pursues accuracy within each frozen subproblem; the second expects it
of the outer sequence
\[
\text{accurate FE model}\ \to\ \text{local approximation}\ \to\
\text{conservative step}\ \to\ \text{new accurate FE model},
\]
no individual frozen approximation being solved to optimality. This is the
stance of sequential convex programming and sequential approximate
optimization, close in spirit to a line search along a predicted direction.
Both readings are legitimate and address genuinely different questions; the
implementation differences follow from the choice of question, not from any
judgement that one viewpoint is superior.

\subsection{First-order accuracy and the eigenvalue regimes}

The data frozen in \eqref{eq:sub}---the matrices $\Krho,\Mrho$, the modes
$\phiv_j$, and the gradients $\fv_{sk}$---all depend on $\rhov$, so the
subproblem is a first-order surrogate whose fidelity degrades with
$\|\Delta\rhov\|$. A distinction must be drawn between the first-order
linearization of the eigenvalue dependence and the convex model actually
optimized: the implementation does not solve \eqref{eq:sub} as a linear
program but optimizes a convex separable MMA approximation, and the two
should not be equated. For a simple eigenvalue the increment enters
\eqref{eq:sub-mult} linearly through $\fv_{nn}$ and carries no intrinsic
curvature from the physics. Absent additional regularization---the move
limit and filtering introduced above---a convex model dominated by such a
curvature-free term has no interior minimizer and is drawn toward the
boundary of the admissible set, i.e.\ the active move limits. This bias
toward large, bound-saturating increments is a tendency of the underlying
linearization, not a literal linear-programming vertex.

This exposes a competition inherent in the inner loop, which we regard as
the central trade-off of the present scheme. Each additional inner MMA
iteration advances the solution of the frozen subproblem toward its optimum,
but it also moves the design away from the point at which the frozen data
were evaluated, and there the approximation deteriorates in two related but
conceptually distinct ways. First, the frozen operators $\Krho,\Mrho$ and
modes $\phiv_j$ increasingly fail to represent the actual design: the finite
element model on which the subproblem rests is no longer that of the current
structure. Second, and independently, the first-order sensitivity model---the
linearization of the eigenvalue dependence through $\fv_{sk}$---loses
predictive accuracy as $\|\Delta\rhov\|$ grows, even relative to the frozen
operators. Refining a fixed approximation that is deteriorating on both
counts improves the solution of the model without necessarily improving the
true objective, and whether continued iteration is worthwhile depends on
which effect prevails. The present implementation rests on the implicit
assumption that, for the problems considered, this joint loss of accuracy
dominates before full convergence of the frozen subproblem becomes
advantageous---so that a single step followed by re-analysis, which restores
a current approximation, is preferable to a converged solve on an outdated
one. We stress that this is a hypothesis supported by the numerical evidence
reported below, not a mathematical proof, and that the balance may differ
for other problem classes or regularizations.

Within the sequential cycle set out above, MMA is a natural choice of step
generator. Its moving asymptotes bound the admissible step and are updated
from the iteration history at the current design; in limiting how far the
variables may move per subproblem they play a role analogous to a
trust-region mechanism, though they are mathematically distinct from it,
bounding and shaping the variables of the convex model rather than
constraining the step within an explicit trust radius. Iterating the frozen
subproblem drives the variables toward those asymptotes, away from the
design at which the model was calibrated.

The force of the two readings is sharpest once the eigenvalue regime is made
explicit, because the frozen subproblem behaves qualitatively differently in
the two cases. In the \emph{simple-eigenvalue} regime---the situation during
most iterations here---the subproblem is the curvature-free first-order
model discussed above: continued optimization of it chiefly drives the
increment toward the model bounds and refines the surrogate rather than the
true objective, which is exactly where the predictor reading is most
compelling and where full convergence of the inner loop stands to gain
least. In the \emph{multiple-eigenvalue} regime the coupling in
\eqref{eq:sub-det} makes the $\Delta\lambda_j$ nonlinear functions of
$\Delta\rhov$, so the subproblem carries genuine curvature that a single
linearized step cannot resolve; iterating it is then far better motivated,
since its accurate solution reflects the true directional behaviour of the
degenerate cluster. We therefore do not conclude that the inner loop is
unnecessary. Rather, its computational importance appears to depend strongly
on the current regime---weakest where the model is essentially first order,
strongest in the multiple case for which it was designed.

The simultaneous constraints on the $J$ lowest modes bear on the same
point. Near a coalescence, several constraints $\beta\le\lambda_j$ may
become active at once, and the associated MMA dual multipliers combine the
corresponding single-mode sensitivities into a step reflecting several
competing modes rather than one. This may play a practical role similar to
that of the generalized-gradient direction that
\eqref{eq:fsk}--\eqref{eq:sub-det} would supply, without forming $\fv_{sk}$
or solving the sub-eigenvalue problem. We regard it as a resemblance in the
computed step, not a demonstrated equivalence: the combination formed from
the MMA duals need not coincide with the direction obtained from the coupled
sub-eigenvalue problem, and we do not claim that it does.

\section{Relation to numerical observations}

The following are consistent with, and motivated by, the behaviour we have
observed on the benchmark configurations.

\emph{Observations.} A literal realization of the nested procedure---the
frozen subproblem iterated to convergence, full increments applied, no
additional damping---was unstable on every benchmark tested: the design
drifted toward near-zero-frequency mechanisms with oscillatory,
non-convergent iterates. The instability was not immediate; it appeared
only after the frozen subproblem had been optimized at length, the
increments then growing large and saturating the box bounds
\eqref{eq:sub-box}, i.e.\ single-step design changes beyond the range in
which the underlying first-order model can be expected to remain accurate.
Rebuilding the approximation after each conservative MMA step, by contrast,
consistently restored stable convergence on the same configurations.

\emph{Interpretation.} These observations are consistent with the analysis
of Section~3: in the simple-eigenvalue regime the curvature-free model
draws its minimizer to the active bounds, and such a step leaves the region
in which the linearization is reliable.

\emph{Hypothesis.} A plausible interpretation is that frequent
re-linearization, rather than convergence of any single frozen subproblem,
is the dominant contributor to stability---as the predictor reading of
Section~3 would suggest. That instability set in only after prolonged
optimization of a fixed approximation, and that rebuilding the model after
each step consistently avoided it, are consistent with the view that the
essential stabilizing mechanism is the return to finite element analysis
before the approximation drifts far from the current design, keeping the
local model current as the optimization proceeds. We regard this as a
working hypothesis supported by the present evidence, not as an established
fact.

\section{Limitations and outlook}

The foregoing does not establish that the algorithm of Du and Olhoff is
mathematically incorrect; it is a convergent scheme whose inner loop is
well motivated for multiple eigenvalues. Our results show only that a
particular literal reconstruction was unstable while a single-step variant
was not, and several alternative explanations cannot be excluded. The
published description leaves undocumented details---most notably the step
control and stopping rule of the inner loop, a sufficiently loose inner
tolerance being, in effect, close to a single-step scheme. Differences in
MMA settings, particularly the initialization and adaptation of the moving
asymptotes, materially affect step size and could account for part of the
gap. Additional stabilization mechanisms may have been present but
unreported. Finally, differences in the finite element model---notably the
interpolation of the low-density element mass, which governs spurious
localized modes and thereby reshapes the eigenvalue landscape near which
\eqref{eq:sub} is built---could contribute independently of the loop
structure.

Taken together, these considerations suggest---as engineering
interpretation rather than proven fact---that the two philosophies are
advantageous under different circumstances. The nested strategy is naturally
favoured when multiple eigenvalues dominate the optimization, so that the
coupled local model carries genuine nonlinear information and solving it
accurately is worthwhile. The present strategy is expected to be favoured
when the optimization proceeds mostly through simple eigenvalues, when
first-order local models lose accuracy rapidly with design change, and when
frequently rebuilding the approximation yields more than repeatedly refining
a single ageing model. On the balance of the available evidence, frequent
rebuilding appears the more important factor for robustness in the regime we
have examined.

These observations raise a broader question that the present work does not
resolve. If the value of solving the local subproblem accurately depends so
strongly on the eigenvalue regime, it is not evident that a sequential
approximation method for eigenfrequency optimization should commit to a
single fixed strategy throughout. The optimal balance between local
optimization and approximation rebuilding may itself be regime-dependent
rather than fixed for the whole run. Whether such an adaptive strategy would
in fact outperform either fixed philosophy is an open question, and a natural
direction for further investigation suggested by the behaviour reported here.

\begin{thebibliography}{9}
\bibitem{DuOlhoff2007} J.~Du and N.~Olhoff.
Topological design of freely vibrating continuum structures for maximum
values of simple and multiple eigenfrequencies and frequency gaps.
\emph{Struct.\ Multidisc.\ Optim.} 34:91--110, 2007.
\bibitem{DuOlhoffErratum} J.~Du and N.~Olhoff. Erratum to
\cite{DuOlhoff2007}. \emph{Struct.\ Multidisc.\ Optim.} 34, 2007.
\bibitem{Svanberg1987} K.~Svanberg. The method of moving asymptotes---a
new method for structural optimization.
\emph{Int.\ J.\ Numer.\ Methods Eng.} 24:359--373, 1987.
\end{thebibliography}

\end{document}
